\paragraph{Item structure.}
Each item in \textsc{AnchorBench-v0.2} specifies:
(i)~a latent signal $\theta \in [15, 85]$,
(ii)~a gold-standard answer $y^{*}$ derived from $\theta$ with
  Gaussian noise ($\sigma{=}12$),
(iii)~low and high anchor values,
(iv)~prompts for each experimental condition, and
(v)~metadata (field, domain, template).
The answer space is constrained to integers in $[0, 100]$,
representing percentages.

\paragraph{Anchor-delivery suites.}
We organize items into five suites:

\begin{itemize}
\item \textbf{External} (30 items): a prompt-embedded anchoring
  question (\eg ``is the rate higher or lower than $X$\%?'')
  appears before the estimation instruction.

\item \textbf{ICL} (12 items): the anchor is delivered via 3-shot
  in-context demonstrations whose answers are biased toward the
  low or high anchor.

\item \textbf{Conversation history} (12 items): the anchor is the
  model's own prior answer (turn 1), re-presented in the context
  of a follow-up question (turn 2).

\item \textbf{RAG} (12 items): anchoring numbers appear in
  retrieved snippets prepended to the prompt.

\item \textbf{Tool} (12 items): anchoring numbers appear in
  simulated tool-call outputs.
\end{itemize}

\noindent Across suites, 30 items belong to the \texttt{external}
suite and 78 to extension suites (ICL, conversation history, RAG,
tool).  Three application fields---health, finance,
operations---contribute 6 domains total.

\paragraph{Conditions.}
Every external-like item has three conditions:
\texttt{control} (no anchor), \texttt{low\_anchor}
(anchor $\leq \theta$), and \texttt{high\_anchor}
(anchor $\geq \theta$).  The anchor gap is 60 percentage points
by construction.
